\documentclass[a4paper]{report}

\usepackage[utf8]{inputenc}
\usepackage{geometry}
\geometry{a4paper, margin=2.5cm}

\begin{document}
	
	% =========================
	%        COPERTINA
	% =========================
	\begin{titlepage}
		\centering
		
		\vspace*{2cm}
		
		{\LARGE Università degli Studi di Trieste\par}
		\vspace{0.5cm}
		{\Large Dipartimento di Ingegneria e Architettura\par}
		
		\vspace{2cm}
		
		{\Huge \textbf{Progetto di Basi di Dati}\par}
		
		\vspace{1cm}
		
		{\Large Sistema informativo per piattaforma sportiva multi-disciplinare\par}
		
		\vspace{2cm}
		
		\begin{flushleft}
			\textbf{Studente:} Enrico Visentin \\
			\textbf{Matricola:} IN0501337 \\
			\textbf{Corso:} Basi di Dati \\
			\textbf{Docente:} 
		\end{flushleft}
		
		\vfill
		
		{\large Anno Accademico 2025/2026\par}
		
	\end{titlepage}
	
	% =========================
	%           INDICE
	% =========================
	\tableofcontents
	\newpage
	
	% =========================
	%       INIZIO DOCUMENTO
	% =========================
	
\chapter{Introduzione}

\section{Introduzione al problema}
Il presente progetto riguarda la progettazione e lo sviluppo di un sistema informativo per la gestione di una piattaforma dedicata ai tornei sportivi multi-disciplina. L’obiettivo del sistema è quello di consentire la registrazione, l’organizzazione e l’analisi di competizioni sportive che coinvolgono diverse discipline, come calcio, basket, pallavolo e tennis.

La piattaforma permette agli utenti di consultare informazioni relative a società sportive, squadre, giocatori, tornei e partite, fornendo inoltre statistiche dettagliate sulle prestazioni individuali e di squadra. Il sistema è pensato per supportare sia la gestione organizzativa dei tornei sia la fruizione dei dati da parte degli utenti finali interessati all’andamento delle competizioni.

\section{Requisiti del problema}
Il sistema deve essere in grado di gestire diverse entità principali, tra cui società sportive, squadre, giocatori, tornei e partite. Ogni società può possedere più squadre, mentre ogni squadra partecipa a uno o più tornei relativi a specifiche discipline sportive.

Per ogni partita devono essere memorizzati i dati relativi alle squadre partecipanti, al punteggio finale, alla data di svolgimento e all’impianto sportivo in cui si è disputata. Inoltre, per ogni giocatore devono essere registrate le informazioni anagrafiche e le statistiche di prestazione relative alle singole partite, come minuti giocati, punti segnati, assist e falli.

Il sistema deve inoltre consentire l’elaborazione di statistiche aggregate, come classifiche, migliori marcatori, prestazioni medie delle squadre e rendimento dei giocatori nei vari tornei.
	
	\chapter{Analisi dei requisiti}
	
	\chapter{Progettazione concettuale (Schema E-R)}
	
	\chapter{Progettazione logica}
	
	\chapter{Schema relazionale}
	
	\chapter{Vincoli e dipendenze funzionali}
	
	\chapter{Implementazione SQL}
	
	\chapter{Query significative}
	
	\chapter{Conclusioni}
	
\end{document}